\documentclass[11pt,a4paper]{article}

\usepackage[utf8]{inputenc}
\usepackage[T1]{fontenc}
\usepackage[french]{babel}
\usepackage[a4paper,margin=2.2cm]{geometry}
\usepackage{amsmath,amssymb,amsthm,mathtools}
\usepackage{lmodern}
\usepackage{enumitem}
\usepackage{hyperref}
\usepackage{xcolor}
\usepackage{fancyhdr}
\usepackage{stmaryrd}
\usepackage{array}
\usepackage{booktabs}
\usepackage{mathrsfs}

\hypersetup{
    colorlinks=true,
    linkcolor=black,
    urlcolor=blue
}

\setlength{\parindent}{0pt}
\setlength{\parskip}{0.6em}

\pagestyle{fancy}
\fancyhf{}
\fancyhead[L]{Corrigé détaillé --- DM Analyse fonctionnelle}
\fancyhead[R]{Terminale générale}
\fancyfoot[C]{\thepage}

\newtheorem*{remarque}{Remarque}
\newtheorem*{idee}{Idée}
\newtheorem*{rappel}{Rappel de cours}

\newcommand{\R}{\mathbb{R}}
\newcommand{\cR}{\mathcal{R}}
\newcommand{\ds}{\displaystyle}
\newcommand{\ppf}{p_f}
\newcommand{\iif}{i_f}
\newcommand{\ppg}{p_g}
\newcommand{\iig}{i_g}

\begin{document}

\begin{center}
    {\LARGE \textbf{Corrigé détaillé --- DM de mathématiques}}\\[0.4em]
    {\large \textbf{Analyse fonctionnelle --- Niveau Terminale}}\\[0.8em]
\end{center}

\hrule
\vspace{0.8em}



\section*{Problème 1 --- Étude d'une fonction définie par ses propriétés locales}

On considère une fonction \(f\) définie sur \(\R\), continue sur \(\R\), dérivable sur \(\R\), vérifiant :
\[
f(0)=0,
\]
\[
\forall x\in\R,\ \forall h\in\R^\ast,\qquad
\frac{f(x+h)-f(x)}{h}\ge 1-\frac{h^2}{1+h^2},
\]
\[
\forall x\in\R,\qquad f(-x)=-f(x).
\]

On peut réécrire l'inégalité de la façon suivante :
\[
1-\frac{h^2}{1+h^2}=\frac{1+h^2-h^2}{1+h^2}=\frac{1}{1+h^2}.
\]
Ainsi, l'hypothèse essentielle devient
\[
\forall x\in\R,\ \forall h\neq 0,\qquad
\frac{f(x+h)-f(x)}{h}\ge \frac{1}{1+h^2}.
\]

\subsection*{Partie A --- Premiers résultats}

\textbf{1. Déterminer la parité de \(f'\), lorsqu'elle existe.}

Comme \(f\) est impaire, on a pour tout \(x\in\R\),
\[
f(-x)=-f(x).
\]

\begin{rappel}
D'après le cours, dérivée d'une composée.
\end{rappel}

On dérive l'identité \(f(-x)=-f(x)\). En utilisant la dérivation d'une composée, on obtient :
\[
-f'(-x)=-f'(x).
\]
Donc
\[
f'(-x)=f'(x).
\]

\[
\boxed{f' \text{ est paire.}}
\]

\bigskip

\textbf{2. Montrer que, pour tout réel \(x\), on a \(f'(x)\ge 1\).}

Fixons un réel \(x\in\R\). Pour tout \(h\neq 0\), on a
\[
\frac{f(x+h)-f(x)}{h}\ge \frac{1}{1+h^2}.
\]

Or, comme \(f\) est dérivable en \(x\), lorsque \(h\to 0\),
\[
\frac{f(x+h)-f(x)}{h}\longrightarrow f'(x).
\]
Et d'autre part,
\[
\frac{1}{1+h^2}\longrightarrow 1.
\]

En passant à la limite dans l'inégalité, on obtient :
\[
f'(x)\ge 1.
\]

Comme \(x\) était quelconque,
\[
\boxed{\forall x\in\R,\quad f'(x)\ge 1.}
\]

\bigskip

\textbf{3. En déduire les variations de \(f\) sur \(\R\).}

\begin{rappel}
Le cours affirme que si \(f'\ge 0\) sur un intervalle, alors \(f\) est croissante sur cet intervalle.
Mieux encore, si \(f'>0\) sur un intervalle, alors \(f\) est strictement croissante.
\end{rappel}

Ici, pour tout \(x\in\R\), on a
\[
f'(x)\ge 1>0.
\]
Donc la dérivée est strictement positive sur \(\R\). Par conséquent,
\[
\boxed{f \text{ est strictement croissante sur } \R.}
\]

\bigskip

\textbf{4. Montrer que, pour tout \(x\ge 0\), on a \(f(x)\ge x\).}

Posons
\[
\varphi(x)=f(x)-x.
\]
La fonction \(\varphi\) est dérivable sur \(\R\), et
\[
\varphi'(x)=f'(x)-1.
\]
Or, d'après la question précédente,
\[
f'(x)\ge 1,
\]
donc
\[
\varphi'(x)\ge 0.
\]
Ainsi, \(\varphi\) est croissante sur \(\R\).

Comme
\[
\varphi(0)=f(0)-0=0,
\]
on déduit que pour tout \(x\ge 0\),
\[
\varphi(x)\ge \varphi(0)=0.
\]
Autrement dit,
\[
f(x)-x\ge 0,
\]
soit
\[
\boxed{\forall x\ge 0,\quad f(x)\ge x.}
\]

\bigskip

\textbf{5. Que peut-on en déduire pour \(x\le 0\) ?}

Soit \(x\le 0\). Alors \(-x\ge 0\). D'après la question 4,
\[
f(-x)\ge -x.
\]
Or \(f\) est impaire, donc
\[
f(-x)=-f(x).
\]
Ainsi,
\[
-f(x)\ge -x.
\]
En multipliant par \(-1\), ce qui renverse le sens de l'inégalité, on obtient :
\[
f(x)\le x.
\]

Donc
\[
\boxed{\forall x\le 0,\quad f(x)\le x.}
\]

\bigskip

\subsection*{Partie B --- Encadrement, convexité intuitive, comparaison de taux de variation}

On suppose dans cette partie que, pour tous réels \(x<y<z\),
\[
\frac{f(y)-f(x)}{y-x}\le \frac{f(z)-f(y)}{z-y}.
\]

\textbf{1. Montrer que pour tous réels \(a<b\),}
\[
f\!\left(\frac{a+b}{2}\right)\le \frac{f(a)+f(b)}{2}.
\]

Posons
\[
m=\frac{a+b}{2}.
\]
Comme \(a<b\), on a bien
\[
a<m<b.
\]
On peut donc appliquer l'hypothèse avec \(x=a\), \(y=m\), \(z=b\), ce qui donne :
\[
\frac{f(m)-f(a)}{m-a}\le \frac{f(b)-f(m)}{b-m}.
\]

Or
\[
m-a=\frac{a+b}{2}-a=\frac{b-a}{2},
\qquad
b-m=b-\frac{a+b}{2}=\frac{b-a}{2}.
\]
Les dénominateurs sont donc égaux et strictement positifs. On en déduit :
\[
f(m)-f(a)\le f(b)-f(m).
\]
Donc
\[
2f(m)\le f(a)+f(b).
\]
Finalement,
\[
\boxed{f\!\left(\frac{a+b}{2}\right)\le \frac{f(a)+f(b)}{2}.}
\]

\bigskip

\textbf{2. Montrer que pour tout réel \(x\),}
\[
f(x)+f(-x)\ge 2f(0).
\]

On applique la question précédente avec \(a=-x\) et \(b=x\). On obtient :
\[
f\!\left(\frac{-x+x}{2}\right)\le \frac{f(-x)+f(x)}{2}.
\]
Comme
\[
\frac{-x+x}{2}=0,
\]
cela donne
\[
f(0)\le \frac{f(-x)+f(x)}{2}.
\]
En multipliant par \(2\),
\[
\boxed{f(x)+f(-x)\ge 2f(0).}
\]

\bigskip

\textbf{3. Confronter ce résultat à la propriété de parité de \(f\).}

Comme \(f\) est impaire,
\[
f(-x)=-f(x).
\]
Donc, pour tout \(x\in\R\),
\[
f(x)+f(-x)=f(x)-f(x)=0.
\]
L'inégalité précédente devient alors
\[
0\ge 2f(0).
\]
Ainsi,
\[
\boxed{f(0)\le 0.}
\]

\bigskip

\textbf{4. Que peut-on conclure sur \(f(0)\) ? Ce résultat était-il déjà connu ?}

On vient d'obtenir
\[
f(0)\le 0.
\]
Mais, dans les hypothèses de départ, on savait déjà que
\[
f(0)=0.
\]
Donc le résultat est compatible avec les hypothèses, mais il n'apporte pas d'information nouvelle.

\[
\boxed{f(0)=0,\ \text{et ce résultat était déjà donné dans l'énoncé.}}
\]

\bigskip

\subsection*{Partie D --- Comportement aux bornes et limites}

On conserve toutes les hypothèses des parties précédentes.

\textbf{1. Montrer que}
\[
\lim_{x\to +\infty} f(x)=+\infty.
\]

D'après la partie A, pour tout \(x\ge 0\),
\[
f(x)\ge x.
\]
Or
\[
x\longrightarrow +\infty \quad \text{quand } x\to +\infty.
\]
Par comparaison, on obtient
\[
f(x)\longrightarrow +\infty.
\]

Ainsi,
\[
\boxed{\lim_{x\to +\infty} f(x)=+\infty.}
\]

\bigskip

\textbf{2. Montrer aussi que}
\[
\lim_{x\to -\infty} f(x)=-\infty.
\]

Soit \(x\to -\infty\). Alors \(-x\to +\infty\). D'après la question précédente,
\[
f(-x)\to +\infty.
\]
Mais \(f\) est impaire, donc
\[
f(x)=-f(-x).
\]
Ainsi,
\[
f(x)\to -\infty.
\]

Donc
\[
\boxed{\lim_{x\to -\infty} f(x)=-\infty.}
\]

\bigskip

On pose maintenant
\[
\varphi(x)=f(x)-x.
\]

\textbf{3. Étudier la parité de \(\varphi\).}

Pour tout \(x\in\R\),
\[
\varphi(-x)=f(-x)-(-x)=f(-x)+x.
\]
Comme \(f\) est impaire,
\[
f(-x)=-f(x),
\]
donc
\[
\varphi(-x)=-f(x)+x=-(f(x)-x)=-\varphi(x).
\]

Ainsi,
\[
\boxed{\varphi \text{ est impaire.}}
\]

\bigskip

\section*{Problème 2 --- Décomposition paire/impaire et structure de l'ensemble des fonctions continues}

On note \(\mathcal C(\R)\) l'ensemble des fonctions continues de \(\R\) dans \(\R\).

Pour toute fonction \(f\in \mathcal C(\R)\), on définit
\[
p_f(x)=\frac{f(x)+f(-x)}{2},
\qquad
i_f(x)=\frac{f(x)-f(-x)}{2}.
\]

\subsection*{Partie A --- Mise en place}

\textbf{1. Montrer que \(p_f\) et \(i_f\) appartiennent à \(\mathcal C(\R)\).}

Comme \(f\) est continue sur \(\R\), la fonction \(x\mapsto f(-x)\) est également continue sur \(\R\), car elle est obtenue par composition de la fonction continue \(f\) avec la fonction \(x\mapsto -x\), elle-même continue.

Ainsi, les fonctions
\[
x\mapsto f(x)+f(-x)
\quad \text{et} \quad
x\mapsto f(x)-f(-x)
\]
sont continues sur \(\R\), en tant que somme et différence de fonctions continues.

En divisant par \(2\), on en déduit que \(p_f\) et \(i_f\) sont continues sur \(\R\).

\[
\boxed{p_f\in \mathcal C(\R)\quad \text{et}\quad i_f\in \mathcal C(\R).}
\]

\bigskip

\textbf{2. Montrer que \(p_f\) est paire et que \(i_f\) est impaire.}

Calculons \(p_f(-x)\) :
\[
p_f(-x)=\frac{f(-x)+f(-(-x))}{2}
=\frac{f(-x)+f(x)}{2}
=\frac{f(x)+f(-x)}{2}
=p_f(x).
\]
Donc \(p_f\) est paire.

Calculons ensuite \(i_f(-x)\) :
\[
i_f(-x)=\frac{f(-x)-f(-(-x))}{2}
=\frac{f(-x)-f(x)}{2}
=-\frac{f(x)-f(-x)}{2}
=-i_f(x).
\]
Donc \(i_f\) est impaire.

\[
\boxed{p_f \text{ est paire et } i_f \text{ est impaire.}}
\]

\bigskip

\textbf{3. Montrer que}
\[
f=p_f+i_f.
\]

Pour tout \(x\in\R\),
\[
p_f(x)+i_f(x)
=\frac{f(x)+f(-x)}{2}+\frac{f(x)-f(-x)}{2}.
\]
En regroupant,
\[
p_f(x)+i_f(x)=\frac{2f(x)}{2}=f(x).
\]
Ainsi, pour tout \(x\),
\[
f(x)=p_f(x)+i_f(x).
\]

Donc
\[
\boxed{f=p_f+i_f.}
\]

\bigskip

\subsection*{Partie B --- Unicité de la décomposition}

\textbf{1. Soient \(u\) une fonction paire et \(v\) une fonction impaire telles que \(f=u+v\). Montrer que \(u=p_f\) et \(v=i_f\).}

On suppose que
\[
f=u+v,
\]
avec \(u\) paire et \(v\) impaire.

Alors, pour tout \(x\in\R\),
\[
f(x)=u(x)+v(x).
\]
De plus,
\[
f(-x)=u(-x)+v(-x).
\]
Comme \(u\) est paire et \(v\) impaire,
\[
u(-x)=u(x),\qquad v(-x)=-v(x).
\]
Ainsi,
\[
f(-x)=u(x)-v(x).
\]

Additionnons les deux égalités :
\[
f(x)+f(-x)=2u(x).
\]
Donc
\[
u(x)=\frac{f(x)+f(-x)}{2}=p_f(x).
\]

Soustrayons maintenant :
\[
f(x)-f(-x)=2v(x).
\]
Donc
\[
v(x)=\frac{f(x)-f(-x)}{2}=i_f(x).
\]

Ainsi,
\[
\boxed{u=p_f \quad \text{et} \quad v=i_f.}
\]

\bigskip

\textbf{2. En déduire que toute fonction continue sur \(\R\) s'écrit d'une et une seule manière comme somme d'une fonction paire et d'une fonction impaire.}

D'après la partie A, toute fonction \(f\in \mathcal C(\R)\) s'écrit
\[
f=p_f+i_f,
\]
avec \(p_f\) paire et \(i_f\) impaire.

Cela prouve \textbf{l'existence} d'une telle décomposition.

La question précédente montre que si \(f=u+v\), avec \(u\) paire et \(v\) impaire, alors nécessairement
\[
u=p_f,\qquad v=i_f.
\]
Il n'y a donc pas deux décompositions différentes possibles.

On a donc montré à la fois l'existence et l'unicité :
\[
\boxed{\text{Toute fonction continue sur }\R\text{ s'écrit d'une et une seule manière comme somme d'une fonction paire et d'une fonction impaire.}}
\]

\bigskip

\textbf{3. Montrer qu'une fonction à la fois paire et impaire est nécessairement nulle.}

Soit \(h\) une fonction à la fois paire et impaire. Alors, pour tout \(x\in\R\),
\[
h(-x)=h(x)
\quad \text{et} \quad
h(-x)=-h(x).
\]
En identifiant les deux expressions,
\[
h(x)=-h(x).
\]
Donc
\[
2h(x)=0,
\]
puis
\[
h(x)=0.
\]
Comme cela vaut pour tout \(x\in\R\), la fonction \(h\) est la fonction nulle.

\[
\boxed{\text{Une fonction à la fois paire et impaire est nécessairement nulle.}}
\]

\bigskip

\textbf{4. Expliquer en quoi ce résultat est lié à l'unicité précédente.}

Supposons qu'une fonction \(f\) admette deux décompositions :
\[
f=u_1+v_1=u_2+v_2,
\]
où \(u_1,u_2\) sont paires et \(v_1,v_2\) impaires.

Alors
\[
u_1-u_2=v_2-v_1.
\]
Le membre de gauche est une fonction paire, car différence de fonctions paires.
Le membre de droite est une fonction impaire, car différence de fonctions impaires.

Ainsi, la fonction
\[
u_1-u_2
\]
est à la fois paire et impaire. D'après la question précédente, elle est donc nulle. Ainsi
\[
u_1-u_2=0 \quad \Longrightarrow \quad u_1=u_2.
\]
Puis
\[
v_1=v_2.
\]

C'est exactement ce qui prouve l'unicité de la décomposition.

\bigskip

\subsection*{Partie C --- Applications analytiques}

On considère désormais une fonction \(f\) continue sur \(\R\).

\textbf{1. Montrer que si \(f\) est dérivable en \(0\) et paire, alors \(f'(0)=0\).}

Comme \(f\) est paire,
\[
f(-x)=f(x)
\qquad \text{pour tout } x.
\]
En particulier, pour tout \(h\),
\[
f(-h)=f(h).
\]

Considérons le taux d'accroissement en \(0\) :
\[
\frac{f(h)-f(0)}{h}.
\]
En remplaçant \(h\) par \(-h\), on obtient
\[
\frac{f(-h)-f(0)}{-h}=\frac{f(h)-f(0)}{-h}=-\frac{f(h)-f(0)}{h}.
\]

Comme \(f\) est dérivable en \(0\), lorsque \(h\to 0\), les deux expressions convergent vers \(f'(0)\).
On obtient donc
\[
f'(0)=-f'(0).
\]
Ainsi,
\[
2f'(0)=0,
\]
d'où
\[
\boxed{f'(0)=0.}
\]

\bigskip

\textbf{2. Montrer que si \(f\) est dérivable et impaire, alors \(f'\) est paire.}

Si \(f\) est impaire, alors
\[
f(-x)=-f(x).
\]
En dérivant cette relation, on obtient :
\[
-f'(-x)=-f'(x).
\]
Donc
\[
f'(-x)=f'(x).
\]
Ainsi,
\[
\boxed{f' \text{ est paire.}}
\]

\bigskip

\textbf{3. Montrer que si \(f\) est dérivable et paire, alors \(f'\) est impaire.}

Si \(f\) est paire, alors
\[
f(-x)=f(x).
\]
En dérivant,
\[
-f'(-x)=f'(x).
\]
Donc
\[
f'(-x)=-f'(x).
\]
Ainsi,
\[
\boxed{f' \text{ est impaire.}}
\]

\bigskip

\textbf{4. Soit \(f\) dérivable sur \(\R\). Exprimer \(p_f'\) et \(i_f'\) à l'aide de \(f'\).}

On rappelle que
\[
p_f(x)=\frac{f(x)+f(-x)}{2},
\qquad
i_f(x)=\frac{f(x)-f(-x)}{2}.
\]

On dérive :
\[
p_f'(x)=\frac{f'(x)-f'(-x)}{2}.
\]
En effet, la dérivée de \(f(-x)\) vaut \(-f'(-x)\).

De même,
\[
i_f'(x)=\frac{f'(x)+f'(-x)}{2}.
\]

Autrement dit,
\[
\boxed{p_f'(x)=\frac{f'(x)-f'(-x)}{2}=i_{f'}(x)}
\]
et
\[
\boxed{i_f'(x)=\frac{f'(x)+f'(-x)}{2}=p_{f'}(x).}
\]

\bigskip

\textbf{5. On suppose maintenant \(f\) deux fois dérivable. Discuter la parité de \(f''\) selon celle de \(f\).}

\begin{itemize}[label=\(\bullet\)]
    \item Si \(f\) est paire, alors d'après la question 3, \(f'\) est impaire. En dérivant encore, la dérivée d'une fonction impaire est paire. Donc \(f''\) est paire.
    
    \item Si \(f\) est impaire, alors d'après la question 2, \(f'\) est paire. En dérivant encore, la dérivée d'une fonction paire est impaire. Donc \(f''\) est impaire.
\end{itemize}

On peut résumer par
\[
\boxed{
\begin{cases}
f \text{ paire } \Longrightarrow f'' \text{ paire},\\[0.3em]
f \text{ impaire } \Longrightarrow f'' \text{ impaire}.
\end{cases}}
\]

\bigskip

\subsection*{Partie D --- Application directe}

On considère la fonction \(g\) définie sur \(\R\) par
\[
g(x)=x^3-3x^2+4x.
\]

\textbf{1. Calculer \(g(-x)\). La fonction \(g\) est-elle paire, impaire, ou ni paire ni impaire ?}

On calcule :
\[
g(-x)=(-x)^3-3(-x)^2+4(-x)=-x^3-3x^2-4x.
\]

Comparons avec \(g(x)\) et \(-g(x)\) :
\[
g(x)=x^3-3x^2+4x,
\]
\[
-g(x)=-x^3+3x^2-4x.
\]

On voit que
\[
g(-x)\neq g(x)
\quad \text{et} \quad
g(-x)\neq -g(x).
\]
Donc
\[
\boxed{g \text{ n'est ni paire ni impaire.}}
\]

\bigskip

\textbf{2. Déterminer la partie paire \(p_g\) et la partie impaire \(i_g\) de \(g\).}

Par définition,
\[
p_g(x)=\frac{g(x)+g(-x)}{2},
\qquad
i_g(x)=\frac{g(x)-g(-x)}{2}.
\]

Calculons d'abord \(p_g\) :
\[
g(x)+g(-x)
=(x^3-3x^2+4x)+(-x^3-3x^2-4x)
=-6x^2.
\]
Donc
\[
p_g(x)=\frac{-6x^2}{2}=-3x^2.
\]

Calculons ensuite \(i_g\) :
\[
g(x)-g(-x)
=(x^3-3x^2+4x)-(-x^3-3x^2-4x)
=2x^3+8x.
\]
Donc
\[
i_g(x)=\frac{2x^3+8x}{2}=x^3+4x.
\]

Ainsi,
\[
\boxed{p_g(x)=-3x^2,\qquad i_g(x)=x^3+4x.}
\]

\bigskip

\textbf{3. Vérifier que \(g=p_g+i_g\).}

On additionne :
\[
p_g(x)+i_g(x)=-3x^2+(x^3+4x)=x^3-3x^2+4x=g(x).
\]
Donc
\[
\boxed{g=p_g+i_g.}
\]

\bigskip

\textbf{4. Montrer que \(p_g\) est paire et que \(i_g\) est impaire.}

On a
\[
p_g(x)=-3x^2.
\]
Alors
\[
p_g(-x)=-3(-x)^2=-3x^2=p_g(x).
\]
Donc \(p_g\) est paire.

D'autre part,
\[
i_g(x)=x^3+4x.
\]
Alors
\[
i_g(-x)=(-x)^3+4(-x)=-x^3-4x=-(x^3+4x)=-i_g(x).
\]
Donc \(i_g\) est impaire.

Ainsi,
\[
\boxed{p_g \text{ est paire et } i_g \text{ est impaire.}}
\]

\bigskip

\textbf{5. Calculer \(g'(x)\).}

Comme
\[
g(x)=x^3-3x^2+4x,
\]
on dérive terme à terme :
\[
g'(x)=3x^2-6x+4.
\]

\[
\boxed{g'(x)=3x^2-6x+4.}
\]

\bigskip

\textbf{6. Étudier le signe de \(g'(x)\).}

On considère le trinôme
\[
g'(x)=3x^2-6x+4.
\]
Calculons son discriminant :
\[
\Delta=(-6)^2-4\times 3\times 4=36-48=-12.
\]
On a donc
\[
\Delta<0.
\]
Comme le coefficient dominant vaut \(3>0\), le trinôme est strictement positif sur \(\R\).

Ainsi,
\[
\boxed{\forall x\in\R,\quad g'(x)>0.}
\]

\bigskip

\textbf{7. En déduire les variations de \(g\) sur \(\R\).}

Comme \(g'(x)>0\) pour tout \(x\in\R\), la fonction \(g\) est strictement croissante sur \(\R\).

\[
\boxed{g \text{ est strictement croissante sur } \R.}
\]

\bigskip

\textbf{8. Dresser le tableau de variation de \(g\).}

Comme \(g\) est un polynôme de degré \(3\) dont le coefficient dominant est positif, on sait que
\[
\lim_{x\to -\infty} g(x)=-\infty,
\qquad
\lim_{x\to +\infty} g(x)=+\infty.
\]

Et comme \(g\) est strictement croissante sur \(\R\), son tableau de variation est :

\[
\begin{array}{c|ccc}
x & -\infty & & +\infty\\
\hline
g'(x) & & + & \\
\hline
g(x) & -\infty & \nearrow & +\infty
\end{array}
\]

\bigskip

\textbf{9. Déterminer les extremums éventuels de \(g\).}

Une fonction strictement croissante sur \(\R\) n'admet ni maximum ni minimum sur \(\R\).

Donc
\[
\boxed{g \text{ n'admet aucun extremum sur } \R.}
\]

\bigskip

\textbf{10. Déterminer les limites de \(g\) en \(+\infty\) et en \(-\infty\).}

Le terme dominant de \(g(x)\) est \(x^3\). Ainsi,
\[
\lim_{x\to +\infty} g(x)=+\infty,
\qquad
\lim_{x\to -\infty} g(x)=-\infty.
\]

\[
\boxed{\lim_{x\to +\infty} g(x)=+\infty \quad \text{et} \quad \lim_{x\to -\infty} g(x)=-\infty.}
\]

\bigskip

\textbf{11. Calculer \(p_g'(x)\) et \(i_g'(x)\).}

Comme
\[
p_g(x)=-3x^2,
\qquad
i_g(x)=x^3+4x,
\]
on obtient
\[
p_g'(x)=-6x,
\qquad
i_g'(x)=3x^2+4.
\]

Ainsi,
\[
\boxed{p_g'(x)=-6x,\qquad i_g'(x)=3x^2+4.}
\]

\bigskip

\textbf{12. Étudier les variations de \(p_g\) sur \([0,+\infty[\).}

Sur \([0,+\infty[\), on a
\[
p_g'(x)=-6x\le 0.
\]
Donc \(p_g\) est décroissante sur \([0,+\infty[\).

\[
\boxed{p_g \text{ est décroissante sur } [0,+\infty[.}
\]

\bigskip

\textbf{13. Étudier les variations de \(i_g\) sur \(\R\).}

Pour tout \(x\in\R\),
\[
i_g'(x)=3x^2+4>0.
\]
Donc \(i_g\) est strictement croissante sur \(\R\).

\[
\boxed{i_g \text{ est strictement croissante sur } \R.}
\]

\bigskip

\textbf{14. Comparer l’allure des courbes représentatives de $g$, $p_g$ et $i_g$.}

On rappelle que
\[
g(x)=x^3-3x^2+4x,\qquad p_g(x)=-3x^2,\qquad i_g(x)=x^3+4x.
\]

On commence par comparer les fonctions deux à deux.

\medskip

\textbf{Comparaison de $g$ et $p_g$.}  
On a
\[
g(x)-p_g(x)=\bigl(x^3-3x^2+4x\bigr)-(-3x^2)=x^3+4x=i_g(x).
\]
Or
\[
i_g(x)=x(x^2+4).
\]
Comme $x^2+4>0$ pour tout réel $x$, le signe de $i_g(x)$ est celui de $x$.

Donc :
\[
\begin{cases}
x<0 \Rightarrow g(x)<p_g(x),\\
x=0 \Rightarrow g(x)=p_g(x),\\
x>0 \Rightarrow g(x)>p_g(x).
\end{cases}
\]

\medskip

\textbf{Comparaison de $g$ et $i_g$.}  
On a
\[
g(x)-i_g(x)=\bigl(x^3-3x^2+4x\bigr)-(x^3+4x)=-3x^2.
\]
Ainsi, pour tout $x\in\mathbb{R}$,
\[
g(x)-i_g(x)\le 0,
\]
avec égalité seulement pour $x=0$.

Donc :
\[
\begin{cases}
x\neq 0 \Rightarrow g(x)<i_g(x),\\
x=0 \Rightarrow g(x)=i_g(x).
\end{cases}
\]

\medskip

\textbf{Comparaison de $i_g$ et $p_g$.}  
On a
\[
i_g(x)-p_g(x)=\bigl(x^3+4x\bigr)-(-3x^2)=x^3+3x^2+4x=g(x).
\]
Or on a montré plus haut que $g$ est strictement croissante sur $\mathbb{R}$ et que
\[
g(0)=0.
\]
On en déduit :
\[
\begin{cases}
x<0 \Rightarrow g(x)<0 \Rightarrow i_g(x)<p_g(x),\\
x=0 \Rightarrow g(x)=0 \Rightarrow i_g(x)=p_g(x),\\
x>0 \Rightarrow g(x)>0 \Rightarrow i_g(x)>p_g(x).
\end{cases}
\]

\medskip

\textbf{Conclusion graphique.}  
Les trois courbes passent par l’origine $O(0,0)$.

\begin{itemize}
    \item Sur $]-\infty,0[$ :
    \[
    g(x)<i_g(x)<p_g(x).
    \]
    Donc la courbe de $g$ est en dessous de celle de $i_g$, elle-même en dessous de celle de $p_g$.

    \item En $x=0$ :
    \[
    g(0)=i_g(0)=p_g(0)=0.
    \]
    Les trois courbes se coupent à l’origine.

    \item Sur $]0,+\infty[$ :
    \[
    p_g(x)<g(x)<i_g(x).
    \]
    Donc la courbe de $p_g$ est en dessous de celle de $g$, elle-même en dessous de celle de $i_g$.
\end{itemize}

Ainsi, la position relative des trois courbes est complètement déterminée par :
\[
\boxed{
\begin{cases}
x<0 : & g(x)<i_g(x)<p_g(x),\\[0.2cm]
x=0 : & g(x)=i_g(x)=p_g(x),\\[0.2cm]
x>0 : & p_g(x)<g(x)<i_g(x).
\end{cases}}
\]

\bigskip

\subsection*{Partie E --- Étude d'une condition fonctionnelle}

On cherche les fonctions continues \(f\) sur \(\R\) telles que, pour tout réel \(x\),
\[
f(x)+f(-x)=2x^2.
\]

\textbf{1. Montrer qu'une telle fonction ne peut pas être impaire.}

Si \(f\) était impaire, alors pour tout \(x\),
\[
f(-x)=-f(x).
\]
Ainsi,
\[
f(x)+f(-x)=f(x)-f(x)=0.
\]
Mais l'équation imposée par l'énoncé est
\[
f(x)+f(-x)=2x^2.
\]
On obtiendrait donc
\[
0=2x^2 \quad \text{pour tout } x,
\]
ce qui est impossible dès que \(x\neq 0\).

Donc
\[
\boxed{\text{une telle fonction ne peut pas être impaire.}}
\]

\bigskip

\textbf{2. Déterminer sa partie paire.}

Par définition,
\[
p_f(x)=\frac{f(x)+f(-x)}{2}.
\]
Or, l'énoncé donne précisément
\[
f(x)+f(-x)=2x^2.
\]
Donc
\[
p_f(x)=\frac{2x^2}{2}=x^2.
\]

Ainsi,
\[
\boxed{p_f(x)=x^2.}
\]

\bigskip

\textbf{3. Montrer que l'ensemble des solutions peut s'écrire à l'aide d'une fonction impaire arbitraire.}

Toute fonction continue \(f\) s'écrit
\[
f=p_f+i_f.
\]
Comme on vient de montrer que
\[
p_f(x)=x^2,
\]
on obtient
\[
f(x)=x^2+i_f(x).
\]
Or \(i_f\) est une fonction impaire continue.

Réciproquement, si \(u\) est une fonction impaire continue et si l'on pose
\[
f(x)=x^2+u(x),
\]
alors
\[
f(-x)=(-x)^2+u(-x)=x^2-u(x),
\]
d'où
\[
f(x)+f(-x)=x^2+u(x)+x^2-u(x)=2x^2.
\]
La condition est donc bien satisfaite.

Ainsi, l'ensemble des solutions est :
\[
\boxed{f(x)=x^2+u(x),\quad \text{où } u \text{ est une fonction continue impaire arbitraire.}}
\]

\bigskip

\textbf{4. Donner trois exemples distincts de telles fonctions.}

Il suffit de choisir trois fonctions impaires continues \(u\).

\begin{itemize}[label=\(\bullet\)]
    \item Si \(u(x)=0\), alors \(f_1(x)=x^2\).
    \item Si \(u(x)=x\), alors \(f_2(x)=x^2+x\).
    \item Si \(u(x)=x^3\), alors \(f_3(x)=x^2+x^3\).
\end{itemize}

Donc, par exemple,
\[
\boxed{f_1(x)=x^2,\qquad f_2(x)=x^2+x,\qquad f_3(x)=x^2+x^3.}
\]

\bigskip

\textbf{5. Parmi elles, déterminer celles qui sont dérivables sur \(\R\).}

Une solution s'écrit
\[
f(x)=x^2+u(x),
\]
où \(u\) est impaire continue.

La fonction \(x\mapsto x^2\) est dérivable sur \(\R\). Ainsi, \(f\) est dérivable sur \(\R\) si et seulement si \(u\) est dérivable sur \(\R\).

Donc les solutions dérivables sont exactement les fonctions de la forme
\[
\boxed{f(x)=x^2+u(x),\quad \text{où } u \text{ est impaire et dérivable sur } \R.}
\]

\bigskip

\textbf{6. Parmi elles, déterminer celles qui sont monotones sur \([0,+\infty[\).}

Une solution s'écrit
\[
f(x)=x^2+u(x),
\]
où \(u\) est une fonction impaire continue.

Il n'existe pas, sous cette seule forme générale, de caractérisation simple de toutes les solutions monotones sur \([0,+\infty[\), car cela dépend du choix de \(u\).

On peut toutefois préciser :
\begin{itemize}[label=\(\bullet\)]
    \item certaines sont monotones, par exemple \(f(x)=x^2\), qui est croissante sur \([0,+\infty[\) ;
    \item d'autres ne le sont pas, par exemple \(f(x)=x^2-x\), car
    \[
    f'(x)=2x-1,
    \]
    qui est négatif sur \(\left[0,\frac12\right[\) et positif sur \(\left]\frac12,+\infty\right[\).
\end{itemize}

Donc on peut dire que les solutions monotones sont celles pour lesquelles la fonction
\[
x\longmapsto x^2+u(x)
\]
est monotone sur \([0,+\infty[\). Si \(u\) est dérivable, cela revient à étudier le signe de
\[
f'(x)=2x+u'(x).
\]

\bigskip

\section*{Exercice bonus --- Taux de variation sans expression explicite}

Soit \(f\) une fonction définie sur \(\R\), et soit \(a\in\R\). On suppose que, pour tout \(h\neq 0\),
\[
\frac{f(a+h)-f(a)}{h}=2a+h+\varepsilon(h),
\]
où \(\varepsilon(h)\to 0\) quand \(h\to 0\).

\textbf{1. Montrer que \(f\) est dérivable en \(a\).}

\begin{rappel}
Par définition, \(f\) est dérivable en \(a\) si le quotient
\[
\frac{f(a+h)-f(a)}{h}
\]
admet une limite finie lorsque \(h\to 0\).
\end{rappel}

Or, par hypothèse,
\[
\frac{f(a+h)-f(a)}{h}=2a+h+\varepsilon(h).
\]
Quand \(h\to 0\),
\[
2a+h+\varepsilon(h)\longrightarrow 2a+0+0=2a.
\]
Le quotient de différence admet donc une limite finie en \(0\). Par définition,
\[
\boxed{f \text{ est dérivable en } a.}
\]

\bigskip

\textbf{2. Déterminer \(f'(a)\).}

Par définition de la dérivée,
\[
f'(a)=\lim_{h\to 0}\frac{f(a+h)-f(a)}{h}.
\]
Or on vient de voir que cette limite vaut \(2a\). Ainsi,
\[
\boxed{f'(a)=2a.}
\]

\bigskip

\textbf{3. Interpréter précisément le rôle du terme \(2a+h+\varepsilon(h)\).}

L'expression
\[
2a+h+\varepsilon(h)
\]
est une écriture précise du taux d'accroissement de \(f\) en \(a\).

\begin{itemize}[label=\(\bullet\)]
    \item Le terme principal est \(2a\) : c'est la valeur vers laquelle tend le taux d'accroissement, donc c'est la dérivée \(f'(a)\).
    \item Le terme \(h\) est un terme correctif qui tend vers \(0\) quand \(h\to 0\).
    \item Le terme \(\varepsilon(h)\) représente un reste, lui aussi négligeable devant \(1\), puisqu'il tend vers \(0\).
\end{itemize}

Ainsi, l'égalité dit que le quotient de différence est \emph{presque égal} à \(2a\) lorsque \(h\) est petit, ce qui justifie la dérivabilité en \(a\).

\bigskip

\textbf{4. Cette information permet-elle de connaître l'expression de \(f\) ? Justifier soigneusement.}

Non, cette information ne permet pas de connaître l'expression complète de \(f\).

En effet, elle décrit seulement le comportement \textbf{local} de \(f\) au voisinage du point \(a\). Elle nous permet de connaître la dérivée de \(f\) en ce point, mais pas l'expression de \(f\) sur tout \(\R\).

Par exemple, si \(a=0\), les fonctions
\[
f_1(x)=x^2
\qquad \text{et} \qquad
f_2(x)=x^2+x^3
\]
ont toutes deux
\[
f_1'(0)=0
\qquad \text{et} \qquad
f_2'(0)=0,
\]
mais ce ne sont pas les mêmes fonctions.

Donc une information sur le seul taux de variation en un point ne suffit pas à reconstituer la fonction entière.

\[
\boxed{\text{Non, on ne peut pas déterminer l'expression de } f \text{ à partir de cette seule donnée locale.}}
\]

\bigskip
\hrule
\bigskip

                        Fin du corrigé 

\end{document}